%% LyX 2.1.4 created this file.  For more info, see http://www.lyx.org/.
%% Do not edit unless you really know what you are doing.
\documentclass[11pt,letterpaper,oneside,english,DIV13]{scrbook}
\usepackage{mathpazo}
\usepackage{helvet}
\renewcommand{\ttdefault}{cmtt}
\usepackage[T1]{fontenc}
\usepackage[latin9]{inputenc}
\pagestyle{headings}
\setcounter{secnumdepth}{3}
\setcounter{tocdepth}{3}
\setlength{\parskip}{\medskipamount}
\setlength{\parindent}{0pt}
\usepackage{color}
\usepackage{babel}
\usepackage{varioref}
\usepackage{textcomp}
\usepackage{makeidx}
\makeindex
\usepackage{graphicx}
\usepackage{nomencl}
% the following is useful when we have the old nomencl.sty package
\providecommand{\printnomenclature}{\printglossary}
\providecommand{\makenomenclature}{\makeglossary}
\makenomenclature
\usepackage[unicode=true,
 bookmarks=true,bookmarksnumbered=false,bookmarksopen=true,bookmarksopenlevel=2,
 breaklinks=true,pdfborder={0 0 0},backref=false,colorlinks=true]
 {hyperref}
\hypersetup{pdftitle={Users Manual},
 pdfauthor={Synercon Technologies, LLC},
 linkcolor=blue,urlcolor=blue,citecolor=blue,anchorcolor=blue}

\makeatletter

%%%%%%%%%%%%%%%%%%%%%%%%%%%%%% LyX specific LaTeX commands.
\pdfpageheight\paperheight
\pdfpagewidth\paperwidth


%%%%%%%%%%%%%%%%%%%%%%%%%%%%%% User specified LaTeX commands.
% \usepackage[pdf]{pstricks}
\usepackage{auto-pst-pdf}
\usepackage{pstricks}
\usepackage{pst-all}
\usepackage{calc}
\newcommand{\degree}{$^{\circ}$}

\AtBeginDocument{
  \def\labelitemiii{\(\circ\)}
  \def\labelitemiv{\(\diamond\)}
}

\makeatother

\begin{document}
\newcommand*{\titleGP}{\begingroup % Create the command for including the title page in the document
\centering % Center all text
\vspace*{\baselineskip} % White space at the top of the page 
\begin{figure}[tbph] 
\centering 
\includegraphics[width=0.15\linewidth]{../media/logos/synercon_logo_v3_only}
\label{fig:logo}
\end{figure}
\vspace{4cm}
{\Huge \textbf{Forensic Link Adapter Field Guide}}\\[2\baselineskip] % Title
{\Huge Synercon Technologies LLC}\\[2\baselineskip]
{\Large For Software Version 1}
\vfill % Whitespace between editor names and publisher logo
\textcopyright {\scshape 2015} \\[0.3\baselineskip] % Year 
Synercon Technologies LLC
\endgroup}
\titleGP


\title{\includegraphics[width=4cm]{media/logos/synercon_logo_v3_only}}


\subtitle{Forensic Link Adapter Users Manual}


\subject{Synercon Technologies, LLC}


\date{For Software Version 1}


\publishers{\href{http://www.synercontechnologies.com}{http://www.synercontechnologies.com}}


\uppertitleback{Disclaimer: The Synercon Technologies does not warrant the completeness
and accuracy of the information produced by its software. While best
efforts by industry experts have been made to generate information
from raw digital data, there may be untested variations, add-on packages,
errors or mistakes due to the large amount of data being processed.
It is the sole responsibility of the user of the Synercon Technologies
systems and Software to determine the fitness of the information provided
by the Software for any particular application and for each individual
case. Synercon Technologies cannot warrant any loss of functionality
due to the lack of reliable network connection to the Synercon Server.
The user is responsible for establishing a reliable Internet connection
for the Software to fully function. Synercon does not warrant that
the Equipment and/or Software will be compatible with each vehicle
model. Electronic components installed on the vehicles may be replacement
parts and not programmed to match the implementations of the Equipment,
hence Synercon cannot warrant the existence or accuracy of the data
recorded on the vehicle or the existence of a fault-free environment.
The information provided by the Software is to be used only by vehicle
manufacturer product experts, engineers, professional accident investigators
and traffic crash reconstruction specialists. The user should agree
that the use of the Software and Equipment requires an appropriate
level of skill and competence. Since Synercon cannot warrant that
the Equipment will create an environment that does not set new fault
codes in the vehicle\textquoteright s electronic control module(s)
which may alter the digital data contained therein, the User assumes
all risks in using the Equipment for the purposes of preserving a
digital record. The user agrees to verify the Equipment is working
according to this manual before using it to support an investigation
or reconstruction. Synercon Technologies, LLC assumes no liability
for incidental, consequential, special or punitive damages, including
but not limited to lost profits and loss of data. The user agrees
to indemnify and hold harmless Synercon Technologies, LLC, any of
its employees, owners, officers, and managers from any legal action
involving the use of the data produced or stored by the software.}

\maketitle
\tableofcontents{}


\chapter{About this Document}

This manual is intended to teach the operator the features and use
of the Forensic Link Adapter (FLA). This is manual is designed as
a teaching tool, rather than a reference guide. There are many acronyms
and terminology used throughout this document, the reader is adviced
to familiarize themselves with the Terminology section before reading.
There are also several other documentation materials provided with
the FLA. The reader is encouraged to examin all documents before using
the FLA to help ensure the most data, and most accurate data is retrieved.

The Field Guide provides an overview for use of the FLA in typical
senarios. This section also provides troubleshooting advice for common
problems.

The ECM Quick References provide a very brief step-by-step guide to
aquiring data from different ECMs. These are intended to be quick
procedures for persons familiar with the FLA and general opperations.


\chapter{Synercon Products}

There are four main components of using Synercon Technologies products.
They will be refered to in the rest of this manual, so the reader
is advised to familiarize themselved with these components to better
understand the rest of this manual.


\section{The Forensic Link Adapter (FLA)}

The FLA is a device that communicates with ECMs using the SAE J1939
and SAE J1587 standards. This device can be a stand-alone data extractor,
or an RP1210 interface to allow the operator to use OEM software.

<picture of FLA>


\section{The FLA Portal}

The FLA Portal is also referenced to as Synercon's Servers. This is
the website \textbackslash{}url\{https://fla.synercontechnologies.com\}.
Data collected from the FLA is uploaded to this site, and operators
are able to view full reports here.

<picture of fla portal>


\section{The FLA Preview}

The FLA Preview is a website hosted by each FLA. It is accessed by
entering the IP address of the FLA into a web browser. This website
provides the operator with the ability to preview data in the field,
without an internet connection.

<picture of fla preview>


\section{The Smart Sensor Simulator (SSS)}

The SSS simulates sensors of a typical vehicle for ECMs. It is used
when the FLA cannot communicate with the vehicle due to damaged wiring.
The SSS reduces the chance the ECM will throw a fault compared to
a standard calibration harness. SSSs are specific to a type of ECM.

<picture of sss>


\chapter{Getting Started}

Congratulations on purchasing the only heavy vehicle scan tool purpose
built for examining vehicles in a forensic or investigative context.
This section provides a guide to getting the FLA System set up for
the first time.


\section{Register as an Operator}

The fla-admin should send you an email to invite you to be an operator
for your FLA on the FLA Portal. The fla-admin email is initiated by
the Organization administrator inviting a new operator. For more details,
please see Chapter XX. It will be neccessary to have a FLA Portal
account to view any data you collect with your FLA.


\section{Opperating the FLA}

The FLA has a 4x20 character LCD screen, along with two buttons to
allow the user to interact with the device.


\subsection{Interacting with the FLA}

There are two general types of screens the FLA will display.

1) <picture of main screen>

The prompt screens will have one or two options on the last line,
corresponding to the button on that side of the FLA.

2) <picture of loading screens>

The working screens will not have an option on the last line, but
usually a percentage bar. During these screens, the FLA is preforming
the operation displayed. \textbackslash{}textbf\{NOTE, during these
screens the RED button can cancle the current task\}.


\subsection{Connect the FLA to the Internet}

Connecting the FLA to the internet allows the device to get updates,
set the system time, and upload data to the FLA Portal.

To connect the FLA to the internet, plug the device into a functional
ethernet port using the included ethernet cable.

<picture of ethernet cable>

The FLA should display an IP address on the second line of the title
screen, the screen that is displayed right after the FLA finishes
booting.

If the FLA displays ``Ethernet Unplugged'' or ``Finding IP Address''
for more than 60 seconds, consult the troubleshooting section.


\subsection{Setting the Time}

Once connected to the internet, the Forensic Link Adapter can synchronoise
with universal time using a dedicated time server. This step is only
necessary if the time displayed on the title screen is not the correct
time.

To correct the FLA time, scroll to the System Configuration screen
on the FLA, using the front-pannel interface.

<picture of sys\_conf screen>

Enter the configuration, then scroll to the Update System Clock screen

<picture of sys\_conf\_update\_clock>

Select ``Update'' with the Green button.

Once the FLA retrieves the current time, it will confirm the new time
with the operator before setting the time.

The FLA will continue to update this time, so the user does not need
to imediatly confirm the time.


\subsection{\label{sub:Setting-the-Timezone-1}Setting the Timezone}

The timezone for the FLA can be set without an internet connection.
To set the system timezone, scroll to the System Configuration screen
(See the above section), then scroll to the Change FLA Timezone screen.
Select ``Next'' to view the next option, and ``Set'' to set the
timezone.

<picture of sys\_conf\_update\_timezone>


\subsection{Updating the FLA}

Updating the FLA provides the latest software, with the most features
and highest chance of getting the most data. It is good practice to
preform an update before downloading a vehicle, or before going out
in the field to download.

To preform an update, folow these steps.
\begin{enumerate}
\item Ensure the FLA is plugged into an ethernet jack, and has an IP address.
\item Scroll to the System Configuration screen, and enter the system configuration
menu <pic of sys\_conf>
\item The FLA will now display the Update FLA Software screen. If the FLA
does not have the Update option, consult the troubleshooting section.
<pic of sys\_conf\_update>
\item Confirm the update <pic of sys\_conf\_update\_confirmation>
\item Wait for the FLA to preform the update. This can take several minutes.
<pif of sys\_conf\_update\_inprogress>
\item When the FLA is finished updating, it will need to shutdown and be
rebooted for the updates to be applied. Select Shutdown. <pic of sys\_conf\_update\_shutdown>
\item After the FLA has completed the shutdown, unplug all sources of power,
If the operator wishes to continue interacting with the device, wait
10 seconds before plugging in power.
\end{enumerate}

\section{Familiarize yourself with the FLA}

After the operator has preformed these steps, the FLA is ready for
field use. The best way to maxamize the validity and completeness
of data collected, the operator is advised to continue reading this
manual. After the user is familiar with this, the best way to get
confortable with using the FLA is to practice. This manual, and the
other supporting documents will help guide operators through the download
and review process.


\chapter{Using the Forensic Link Adapter}


\section{Specifications}

The Forensic Link Adapter is based on a powerful processor based on
the Texas Instruments AM335x 1GHz ARM\textregistered{} Cortex-A8 Processor.
It has 512MB of DDR3 RAM, 16 GB of internal storage, the NEON floating-point
accelerator and 2 PRU 32-bit microcontrollers. The device has 2 Controller
Area Network (CAN) channels, 2 J1708 Serial Channels, 100/10 Mbps
Ethernet, USB 2.0, and a full implementation of the DG tech DPA4+
for RP1210 Compliance. The device has a GPS receiver, Real Time Clock,
an accelerometer, and a rate Gyro. The display uses Organic LEDs for
better sunlight readability. It runs embedded Linux and serves its
own website.


\section{Working with the Forensic Link Adapter}

This section will explain the basics of using the FLA. This is intended
for new users.


\subsection{Powering Up}

The FLA will begin to boot whenever a power is supplied to the device.
This can be done in one of three ways
\begin{enumerate}
\item Using the vehicle connection cable to power the FLA using the battery
of the vehicle
\item Using the wall adapter to power the FLA using a wall outlet
\item Using the 12v cigarrette adapter to power the FLA either from a vehicle
battery, or a portable battery pack
\end{enumerate}
Once the FLA is finished booting, the device will display the title
screen <pic of title\_screen>. When this screen appears, the device
is ready to opperate. If the device does not reach this screen in
60 seconds, see the troubleshooting section of this manual.

Once an investigator has established permission to download the data
from a vehicle, the key needs to be turned to the on (not start) position
for the FLA to communicate with the Electronic Control Modules (ECMs).
The FLA can be plugged into the 9-pin diagnostic connector, which
is usually located below the dash on the left side or to the left
of the driver's seat. If the ECM was removed, then it should be connected
through a Smart Sensor Simulator that emulates a truck. 


\subsection{Shutting Down}

The FLA needs to be shutdown correctly in order to minimize the chance
of loosing data. To shutdown the FLA properly, follow these steps.
\begin{enumerate}
\item Scroll to the Shutdown screen, and select shutdown <pic of shutdown\_screen>
\item Confirm the shutdown <pic of shutdown\_confirm\_screen>
\item Wait until the FLA indicates it is safe to remove power <pic of safe\_to\_unplug>
\end{enumerate}

\section{Display Screen Menu System}

The FLA uses the front pannel display as the primary way to interact
with the unit. This section discribes the design of the menu system
and how to navigate it.


\section{Menu Screen Descriptions\label{sec:Menu-Screen-Descriptions}}

The menu system consists of the main menu, and a Systems Configuration
menu.

<Drawing of menu>
\begin{enumerate}
\item \label{enu:Startup}A Startup Screen is shown when booting the device.
A progress bar scrolls across the bottom in about 20 seconds. After
that, the progress bar may change to say the words ``Please Wait...''
If the message to ``Please Wait...'' stays on and the automatic
transition to display \ref{enu:Execute} does not happen within 30
seconds, the system may not have booted. The FLA may need to be unplugged
and restarted if this screen does not automatically go away. See Section
\vref{sec:FailsToBoot} for further troubleshooting guidance.
\item \label{enu:Execute}The Program Loading screen shows the user that
the FLA booted successfully and is loading the program needed to interface
with the user through the display. It also sets up the files and directories
for the current session. A few things happen in the background while
this screen is present:

\begin{enumerate}
\item The FLA system time is updated. If an Internet connection is available
(i.e. the Ethernet cable is plugged into a known good network), then
the FLA will get its system time from the network time server at NIST.
Once it automatically updates the time, it also updates the time on
its battery powered real time clock (RTC). If an Internet connection
is not available, then the system time is updated from the RTC. 
\item The FLA tries to obtain a GPS lock.
\item The FLA looks to see if vehicle networks are present. It detects the
presence of network traffic on the J1708 network (9600 baud) and the
J1939 network (CAN bus).
\end{enumerate}
\item \label{enu:Status}The first line of the Status Screen shows the current
time with the time zone. The system time zone can be changed from
the web interface, as described in Section \vref{sub:Setting-the-Timezone}.
The second line of the status screen shows the status of the Ethernet
connection, as described in Section \vref{sub:IP-Addresses}. On the
third line, the status of the vehicle network connection is displayed.
The following options are available:

\begin{enumerate}
\item \label{enu:StatusA}No Vehicle Net means that network message traffic
was not detected on either the J1939 network (CAN bus) or the J1708/J1587
serial bus. A message will also suggest to ``Check Ing. Key'' to
remind the operator that the ignition key switch must be turned to
the on or run position for network traffic to be present. If the key
is on, sufficient power is available, and no vehicle networks are
present, then the vehicle network may be compromised.
\item \label{enu:StatusB}CAN Bus Only means that messages using the Controller
Area Network are present. Since SAE J1939 specifies CAN as the network,
this detects the presence of J1939 messages. J1708 network traffic
was not detected.
\item \label{enu:StatusC}1708/1587 means no CAN or J1939 messages are present
on the vehicle networks. This may be the case for older vehicles.
\item \label{enu:StatusD}CAN and J1708 mean both networks are present and
have traffic.
\end{enumerate}

The system voltage is also displayed as detected by a voltage divider
connected to the main processor. This voltage indication is not calibrated,
but should give the user some feedback on the health of the vehicle
system voltage. On the bottom line, there are button commands and
an indicator showing the number of GPS satellites. The Red button
advances the FLA to the next screen. The Green button starts the scan
of the vehicle network to identify the system. The scan button is
not available if no vehicle network is present. 

\item \label{enu:Shutdown}The Shutdown dialog enables the user to shutdown
the FLA computer system in a friendly way. This is similar to pressing
the shutdown command on a PC. The shutdown command can also be invoked
from within the FLA web interface. If the operator does not want to
shutdown the FLA, then he or she can press the Red button to scroll
to the next menu option. If the Green button is pushed, then the operator
will invoke the shutdown sequence as described in the following displays

\begin{enumerate}
\item \label{enu:Shutdown-Confirm}A shutdown confirmation asks the user
if that was their intention. This helps catch unintended shutdowns.
The Red button will escape the shutdown command and move to the next
menu option. The Green button will confirm the shutdown operation.
A reminder that all the files will be saved is displayed to the user
as well.
\item \label{enu:Shutdown-Sequence}The shutdown sequence requires power,
so the user is advised not to unplug the device. During this phase,
an LED indicator on the side may turn yellow. 
\item \label{enu:Shutdown-Complete}When shutdown is complete the main processor
has turned off power. The display driver remains on with a static
message that everything is ready to unplug. No commands are available.
The LED should turn Red.
\end{enumerate}
\item \label{enu:Permission}Once the FLA detects a vehicle network and
the user presses the button to scan, the FLA asks the operator if
they have the proper permission to download digital data. This permission
is dependent on the case and the jurisdiction. The FLA system time
is recorded when permission was acknowledged. Pressing the Green button
in this menu defines the start of a download. The next screen to display
is number \ref{enu:Identifying-Vehicle}.
\item \label{enu:Upload-to-Server}The screen that says Upload data on FLA
on the first line enables the user to push data to the server for
archiving, decoding, decrypting and display. Upon entry into this
screen and a valid IP address, the FLA tries to ping the Synercon
Server for 5 seconds or less. After the ping, the FLA will display
either \ref{enu:UploadA} or \ref{enu:UploadB}.

\begin{enumerate}
\item \label{enu:UploadA}If there is a valid network connection to the
Synercon Server, the message on the second line of the FLA display
will say Internet Ready and the Upload option is enabled for the Green
button. Pressing the Green button will display screen \ref{enu:UploadToServer}.
\item \label{enu:UploadB}If there is not a connection to the Network server,
then a message on the second line will say No Network to Server. Be
sure the Ethernet cable is connected, a valid IP address is displayed,
and the network administrator has granted your device access to the
internet. The operator may want to try connecting the FLA to a different
network.
\item \label{enu:UploadToServer}The Upload to Server screen displays the
connection status as ``OK'' or ``Failed.'' When the connection
is valid, the FLA securely uploads data package files to the Synercon
Server. The number of files to upload (NN) corresponds to the number
of times the FLA downloaded a vehicle or ECM since the last time data
packages were sent to the Synercon Server. The number of files uploaded
(MM) shows a counter that increments after each successful upload
to the server. 

\begin{enumerate}
\item \label{enu:Upload-Failed}Upload Failed. An error code is displayed
in place of \texttt{abcdefghijklmno}. This error code will be important
to the Synercon Support staff and may be related to FLA authentication.
\item \label{enu:Upload-Success}Upload Successful. See results at \href{http://fla.synercontechnologies.com/}{http://fla.synercontechnologies.com/}.
This website is where the data will be displayed.
\end{enumerate}
\end{enumerate}
\item \label{enu:Change-IP-network}Change IP network settings. This screen
enables the operator to turn the FLA into a Dynamic Host Configuration
Protocol (DHCP) server and give IP addresses to other computers on
the network. Since only one DHCP server should be on a network at
a time, this feature off by default. However, in the field, the FLA
may need to set up an FLA network so the operator can connect to the
FLA web page. The following descriptions explain how to enable and
disable the DHCP server.

\begin{enumerate}
\item \label{enu:NetworkA}By default the DHCP server is disabled, as shown
on line 2. Pressing the Green button will pull up the confirmation
screen shown as display \ref{enu:DHCP}. If the words ``No Valid
IP Found'' are displayed on line 3, then turning on DHCP may be a
good idea. If a valid IP address (IP: XXX.XXX.XXX.XXX) is displayed,
then DHCP may not be needed. 
\item \label{enu:NetworkB}After the DHCP server has been enabled, line
2 will display as such. The IP address should reset to 10.0.0.1. 
\item \label{enu:DHCP}The field use only warning should remind the operator
that an already established network does not need another DHCP server.
Enabling DHCP on the FLA on an established network may lead to IP
address conflicts for other devices on the network. 
\end{enumerate}
\item \label{enu:GPS-Status}{[}Under development{]} A GPS utility screen
exists to enable the user to see and save GPS information. The GPS
coordinates in decimal are displayed along with speed and direction.

\item It is recommended to update the FLA software periodically. The FLA
must be connected to the Internet in order to get software updates.

\begin{enumerate}
\item \label{enu:UpdateA}If an Internet connection to the update server
exists, then the Green button will be enabled for the user to select
Update. 

\begin{enumerate}
\item \label{enu:UpdateConfirm}Once the Green button is pushed, an update
confirmation display enables the user to cancel an update operation.
Pressing the Green button again will confirm the update should start.
It is important to not unplug the FLA power and Ethernet cables during
the update.
\item \label{enu:Update-Progress}Update progress is displayed with current
operations shown as xxxxxxxxx on display. This process takes a long
time and the scroll bar may seem to pause occasionally.
\end{enumerate}
\item \label{enu:UpdateB}If there is no access to the update server, then
the Green button is disabled. The FLA will need to be connected to
the Internet by the Ethernet cable.
\end{enumerate}
\item \label{enu:Identifying-Vehicle}Once permission was acknowledged,
the FLA will scan the available vehicle networks. The program is looking
for all data defined in the J1587 or J1939 standards that may be available
and useful, as explained in further detail in Section \vref{sec:Standards-Based-Data}.
As those data elements are obtained from their respective networks,
the counter will increment reflecting the recording of the data. Once
the data elements have been requested, the FLA will be able to identify
the different ECMs.
\item \label{enu:Id-Vehicle}The ECM Identification screen shows the identification
information stored in the first Engine ECM. The ECM make (AAAAA) and
model (MMMMMMMM) comes from requesting the Component Identification.
The Component Identification is a J1587 message with PID of 241 or
a J1939 message with SPNs 586 and 587. The Vehicle Identification
Number (VIN) comes from J1587 PID 237 or J1939 SPN 237. These data
may or may not be present depending on how the ECM is programmed.
If the data are not present from the Engine ECM, then the next source
is used, if available. 
\item \label{enu:ECM-Specific-Data}Once the ECM type is known, the FLA
will check to see if the ECM is supported. If it is supported, then
ECM specific data can be obtained. 

\begin{enumerate}
\item \label{enu:ECMa}If the ECM is supported, then different types of
data can be obtained that include Event Data, Configuration Data,
Historical Data, and Fault Data. If this menu screen is displayed,
then the FLA detected that it can capture data from an ECM on the
network. Pressing the Green button will start the process of gathering
specific data.
\item \label{enu:ECMb}No current support for getting ECM Specific Data
exists of this screen is displayed. As such, both the Red and Green
buttons move the interface to the RP1210 mode where the operator can
use the FLA with other software. 
\item \label{enu:GettingECMdata}Getting ECM data... is a display screen
that shows the progress of the current download. The tag <Data Sought>
is dependent on the ECM and may say things like Hard Brake, Engine
Usage Log, Data Plate, or other phrases that to indicate the progress
of the download. After each attempt, a flag saying either Success
or Failed will be displayed. 
\item \label{enu:DoneECMdata}Once the ECM specific data is downloaded,
a message saying Finished ECM Data is displayed along with the total
number of bytes downloaded. If the operator wants to continue an investigation
using PC based software, then he or she can press the Red button to
go into RP1210 passthrough mode. Otherwise, the Green button brings
the user back to the upload menu. In the field. it may be that a data
upload is possible due to the lack of an internet connection. In this
case, the operator can press the Red button once to get to menu \vref{enu:Status}
and again to get to the Shutdown menu.
\end{enumerate}
\item The RP1210 passthrough mode enables operators to use other PC based
software to interface with the vehicle. The FLA will be connected
to a PC by a USB cable to communicate on the vehicle bus. The FLA
uses a DGTech eDPA to communicate with the PC. As such, it is important
to have the DG Drivers installed before using the RP1210 mode, as
described in Section \vref{sub:RP1210-Drivers}. This mode invokes
a write blocking feature that prevents certain messages from reaching
the ECM. For more details on the messages that get blocked when the
write-blocker is running, see Section \vref{sec:Forensic-Write-Blocking}. 

\begin{enumerate}
\item \label{enu:RP1210}This display screen asks to enable the passthrough
mode. If the Red button is pressed, then the system goes back to the
status display (Screen \ref{enu:Status}). Pressing the Green button
turns on the Passthrough mode.
\item \label{enu:RP1210Enabled}Passthrough mode for CAN (J1939) is on when
this screen is displayed. A reminder to plug in the USB cable is printed
on line 2. At any time, the operator can terminate the pasthrough
mode and switch to J1708 pass mode or transition to uploading the
data. This mode is used for all newer ECMs.
\item \label{enu:J1708-pass-mode.}J1708 pass mode is on when this screen
is displayed. A reminder to plug in the USB cable is printed on line
2. At any time, the operator can terminate the pasthrough mode and
switch to J1708 pass mode or transition to uploading the data. This
mode is used for all Caterpillar, DDEC IV and DDEC V ECMs.
\end{enumerate}
\end{enumerate}
Network traffic is recorded and logged during screens \ref{enu:Id-Vehicle},
\ref{enu:GettingECMdata}, and \ref{enu:RP1210Enabled}. A data package
is defined from network data coming from the vehicle starting when
the user acknowledges permission to when the operation of the FLA
transitions back to the Upload option. 


\section{\label{sub:IP-Addresses}IP Addresses}

The IP address Line can have the following statements:
\begin{description}
\item [{IP:~XXX.XXX.XXX.XXX}] A valid IP address. XXX can be any number
between 0 and 255. Often a private home network will have a router
that serves an IP address with the the first numbers of 192.168.XXX.XXX.
\item [{IP:~10.0.0.1}] The IP address when the FLA's DHCP server is activated.
The FLA will give other devices on the network an IP address too,
which may cause problems on an established network. 
\item [{Ethernet~Unplugged}] There is no live connection between the FLA
and an outside device. If a network connection was present, there
may be around a 30 second lag for this message to display after a
cable is unplugged.
\item [{Finding~IP~Address}] The Ethernet connection is plugged in and
linked to another connection, but the FLA has not been issued an IP
address from the network router.
\end{description}



\chapter{The FLA Preview}


\section{Accessing the FLA Preview}

Accessing the FLA Preview is done by entering the IP address of the
FLA into a web browser. Please see the requirements section for more
information about supported web browsers. As shown below, enter the
IP Address exactly as it appears on the IP Address line.

<Pic of correct/ incorrect ip addresses>

This will load the FLA Preview main page, shown below.

<pic of fla\_preview\_main\_page>

This page provides an overview of the current status of the device.
The At a glance section will display the number of pending data packages,
that is the number of data packages that have been downloaded to the
FLA but not yet uploaded to the FLA Portal, and if there is a download
in progress.


\section{Using the FLA Preview}

The FLA Preview has many uses, including allowing operators to preview
data in the field, download the User Manuals, download RP1210 drivers,
and many other useful tasks we will cover in this section.


\subsection{View Local Data}

Under the View Local Data page, the operator is able to view data
on the FLA. There are two sections to local data, Pending Uploads,
and Archived Data. Pending Uploads will show any data packages that
are ready to be uploaded to the FLA Portal.

<pic of local\_data>

This page will give an overview of each pending data package. The
OEM data column will be YES if the FLA was able to recover some OEM
data, and NO if the FLA could not. The yellow Archive button will
remove the data package from being uploaded, and move it to the archived
section. If an archive is unintentionally archived, the FLA will have
to be serviced to retrieve the data. For more information about servicing
the FLA, see the Service section.

When a download is in progress, the row for that entry will indicate
the download is in progress, and will not provide the ``View Data''
button. Once the user finishes the download, the operator can refresh
the page, and will see the ``View Data'' button. For more information
about when a download is finished, and how to preform a download,
refer to the Field Guide.

Selecting the ``View Data'' will provide the user with a preview
of the report. This preview includes all of the standards-based information,
the components found, and any identification information, as well
as any in-field non-standards-based previews. Non-standards-based
information is any data that is provided by OEM software, such as
DDEC reports' XTR files, and CAT Snapshots. The operator will get
a pireview of this data here in accordance with the type of ECM being
analyzed. The full reports will be avaliable through the FLA Portal
when the operator uploads the data. Typically, the following will
be avaliable:

1. Preview of DDEC ECMs

If non-standards-based data is collected from a DDEC ECM, an XTR file
will be avaliable for download.

2. Preview of CAT ECMs

If non-standars-based data is collected from a CAT ECM, a CSV file
will be avaliable containing 5 frames around the trigger for any quick
stop records.

The preview will also contain a SHA 256 Sums file for the operator
to download. This file contains the SHA 256 sums of all data files
the FLA has created. The operator can download this file, and confirm
the data uploaded to the FLA Portal was the same data the FLA created.
For more information on SHA 256 sums, as well as recommended use,
see the SHA 256 sums of suggested practices.

The operator can also click on the ``Archived Data'' button on the
View Local Data page to view data packages that were archived. Data
packages that are uploaded to the FLA Portal are automatically archived.
From here the user can access all of the same downloads and information
as when the data package was pending uploading. Meaning, once the
data is uploaded, the user can still downloa XTR, CSV, or SHA 256
Sum files, as well as any standards-based data aquired. If the FLA
has a large number of archived data packages, the archived page may
take several moments to load.

<pic of archived\_datas>


\chapter{Using the FLA Portal}

The FLA Portal is where operators can view full reports of the data
they collected in the field. Administrators will also use the FLA
Portal to manage FLAs, and users.


\section{Accessing the FLA Portal}

The FLA Portal is accessed by navigating to \href{https://fla.synercontechnologies.com/}{https://fla.synercontechnologies.com/}
in a web browser. For a list of supported web browsers, refer to the
XXX section.

The FLA Portal will require the operator to login, using the username
and password used when the operator created the account. Note that
BOTH username and password are case-sensative.

\includegraphics[width=0.8\textwidth]{\string"media/fla_portal_screenshots/FLA Portal Login Screen\string".png}

Once logged in, the operator will be able to view data aquired from
the assigned FLA.


\section{Using the FLA Portal}

For operators, the main use of the FLA Portal is to view full reports.
Administrators should refer to the Administrative Guide for further
information about administrating on the FLA Portal.

Once an operator uploaded a data package to the FLA Portal, the data
will be avaliable for viewing by everyone in the group. A Logged-in
operator will be able to view all data packages uploaded to the server.

<picture of logged-in operator>

Selecting one of these reports will open the full report.

<picture of full report>

The FLA full report has many sections, including


\subsection{Vehicle Information}

This section contains any components that responded to a component
Identification request. The VIN, Vehicle Odometer, and Engine Hours
will also be displayed under Engine \#1 for the appropriate network.
This section also contains the time records, including vehicle time
if avaliable, and time of the last GPS Lock.

<pic of vehicle information>


\subsection{Data sections}

Data sections, such as Historical Data, Event Data, etc., are sections
that contain either standards-based data, or non-standards-based data.
These sections contain all of the data the FLA has decoded. Due to
the wide variety of ECMs, and ECM configurations, these sections may
contain different data depending on the ECM being analyzed.


\subsection{Nomenclature}

This section contains a breif overview of the nomenclature used in
the report.


\subsection{Network Logs}

This section contains the logs of all of the FLA's interactions with
the ECMs. These logs are provided as a means to audit the messages
being sent. If the RP1210 passthrough is engaged, separate logs for
any passthrough transactions will be avaliable.

<picture of netowrk logs>


\subsection{Version Information}

This section provides the software versions that were used when the
report was generated.

<picture of version information>


\subsection{SHA 256 Sums}

This section provides the SHA 256 Sums for all of the data files the
FLA generated when the download was preformed. These sums are encryped
and sent to the FLA Portal for verification. The FLA Portal will generate
its own SHA 256 Sums of the data files uploaded by the FLA. The FLA
Portal will then compare the two, and verify the integrety of the
upload.

<picture of sha\_256\_verified>

If there is a SHA 256 Mismatch, this may suggest the data has been
tampered with.

<picture of sha\_256\_mismatch>


\subsection{Access Information}

This section provides an audit of who is currently accessing the report,
and the device's IP address.


\chapter{Administrator's Guide}

This section is intended to teach administrators how to manage FLAs
on the FLA Portal.


\section{FLA Portal terminology}

The FLA Portal has the following categories:
\begin{description}
\item [{Organizations}] are the top level with a site administrator as
the person in charge. Typically this is a highway patrol, a state
police, or a company. The person in charge of the organization has
the ability to oversee the data and use of the different FLAs.
\item [{Groups}] are a sub-level within an organization. For example, East
District and West District may be groups within an State Police organization.
If an organization only has a few operators and a few FLAs, it may
not make sense to have different groups. 
\item [{Operators}] are trained users of the FLA and are responsible for
collecting data in the field. Operators can be assigned to different
groups. Every operator in the group is able to see the data packages
from all the FLAs assigned to the group, regardless of which operators
downloaded the data.
\item [{Forensic~Link~Adapters}] are the hardware devices that an organization
owns.
\end{description}

\section{Account Overview}

Forensic Link Adapters are assigned to each organization at the time
of purchase. Once the organization completes the purchase, the serial
numbers of each FLA will be loaded into the Organization's profile.
An invitation to set up the Organization will be sent to the site
administrator. For example, the Lieutenant in charge of the crash
reconstruction division will become the site administrator once he
or she registers by following the link sent in the invitation e-mail.

\begin{figure}
\includegraphics[width=1\textwidth]{media/fla_portal_screenshots/FLAPortalAccountOverviewScreenShot}

\caption{\label{fig:Example-FLA-Account-overview}Example Forensic Link Adapter
Portal web page for the Account overview}
\end{figure}


Once logged in as a Site Administrator, an account overview page will
be available, much like the one shown in Figure \vref{fig:Example-FLA-Account-overview}.
Many of the fields on the web page are linked to various functions.
The top bar (in black) has links to your default login page, the latest
FLA download if you are an operator, a list of all the organizations,
FLAs, and some account access settings. The web page is set up in
three distinct tables as seen in Figure \ref{fig:Example-FLA-Account-overview}.
\begin{enumerate}
\item Organization Groups
\item FLA Operators
\item Forensic Link Adapter (FLA)
\end{enumerate}

\section{Group Management}

It is the responsibility of the Organization administrator to set
up groups and operators according to the desired structure of the
Organization. You can have operators that have access to more than
one group.

\begin{figure}
\includegraphics[width=1\textwidth]{media/fla_portal_screenshots/FLAPortalFieldOperatorsScreenShot}

\caption{\label{fig:Example-Group-Details}Example Group Details web page from
the Forensic Link Adapter Portal.}
\end{figure}



\subsection{Adding an Operator}

There are two ways to get an operator access to the group's FLAs.
The first is to invite a new operator and the second is to invite
an existing operator. However, the existing operator must have accepted
an invitation from the Organization under a unique e-mail address.
By clicking on the lower left box the says \textsf{Invite new operator}
shown in \vref{fig:Example-Group-Details}, a web page similar to
the one shown in \vref{fig:New-Operator-invitation}. The required
user information includes name and email address. The FLA portal uses
e-mail addresses to identify operators. Therefore, if a person belongs
to different organizations, they would have to use a different e-mail
for each organization. New operators will get an e-mail from \texttt{fla-admin@synercontechnologies.com}
with a link to register with the site. Many times this e-mail will
be filtered as junk, so advising the recipient to check their e-mail
filters and junk folder may be necessary.

The Work Address details on the Invite and operator page are used
to fill in a report with 

\begin{figure}
\includegraphics[height=10cm]{media/fla_portal_screenshots/FLAPortalInviteNewOperatorScreenShot}

\caption{\label{fig:New-Operator-invitation}New Operator invitation form on
the FLA portal website.}
\end{figure}



\subsection{Assigning an FLA to the Group}

Each group needs to have at least one Forensic Link Adapter assigned
to it. To assign an FLA to a group, click the button to see a list
of available FLAs for your organization. Only people in the group
can see the data on the FLA assigned to that group

\begin{figure}
\includegraphics[width=1\textwidth]{media/fla_portal_screenshots/FLAPortalAssignedFLASuccessScreenShot}

\caption{Assigned FLAs}
\end{figure}



\chapter{Obtaining Data With Manufacturer's Software}

There are two protocols used to communicate over a heavy truck: 1)
CAN/J1939, and 2) J1708. These modes are needed to make use of the
DG RP1210 embedded device used to communicate with Manufacturer\textquoteright s
software. 


\section{Installing Programs and Drivers}


\subsection{\label{sub:RP1210-Drivers}RP1210 Drivers}

The FLA uses the DG Technologies drivers to enable Windows programs
to communicate with the ECMs. Download and install the drivers on
your laptop or local machine from

\href{http://www.dgtech.com/product/dpa/software/DPA4P_136.zip}{http://www.dgtech.com/product/dpa/software/DPA4P\_{}136.zip}.

For more information, see the DG website

\href{http://www.dgtech.com/product/dpa4plus/downloads/downloads.php}{http://www.dgtech.com/product/dpa4plus/downloads/downloads.php}.
The Adapter Validation Tool is a useful program to test the connections
between the computer, FLA, and Vehicle.


\subsection{Manufacturer's Software}

Download Cummins PowerSpec from

\href{http://cumminsengines.com/powerspec}{http://cumminsengines.com/powerspec}. 

This will only work on Windows 7 or newer. You need to register PowerSpec
for free, but you do not have a license the program for it to work.

Download Detroit Diesel DDEC Reports from

\href{ http://www.ddcsn.com/cps/rde/xchg/ddcsn/hs/3448.htm}{ http://www.ddcsn.com/cps/rde/xchg/ddcsn/hs/3448.htm},
as shown in \vref{fig:Download-page-forDDEC}.

The direct download link is

\href{http://www.ddcsn.com/cps/rde/xbcr/ddcsn/DDECReports805.exe}{http://www.ddcsn.com/cps/rde/xbcr/ddcsn/DDECReports805.exe}.

Restart your computer after installing. 


\subsection{PDF Printer}

If a PDF generator is not installed on your computer system, then
having the ability to print to PDF is useful. An example of a free
PDF printer is the 

\href{http://www.bullzip.com/products/pdf/download.php}{BullZip PDF printer.}


\section{Enabling RP1210 Passthrough Mode}

If the ECM is a Caterpillar, DDEC IV DDEC V, then enable the J1708
Pass-through mode from the FLA menu screen. This is screen \ref{enu:J1708-pass-mode.}.

If the ECM is a Cummins, DDEC VI, DDEC 10, or other then enable CAN/J1939.

Once enabled, connect a USB cable from the FLA to the laptop with
the DG drivers installed.


\section{Cummins PowerSpec Download Protocol}
\begin{enumerate}
\item Turn the ignition key to the on position (if it is not already on),
but do not start the engine.
\item Plug in the FLA to the diagnostic connector. Ensure it powers on and
boots.
\item Perform a Standards based download using FLA Diagnostics.
\item Enable J1939 Passthrough mode
\item Launch DG Adapter Validation Tool (AVT). 

\begin{enumerate}
\item Select the appropriate installed device driver. 
\item Switch the protocol to J1939 in the adapter validation tool.
\item Click Run Test.
\item If the Adapter Validation Tool passes the test, then two windows will
turn green as shown in Figure \vref{fig:Adapter-Validation-Tool}.
This means the ECM and the computer are connected through the RP1210
Device. If a test fails, review the suggestions output by the Adapter
Validation Tool and try again. It may require shutting down the FLA,
unplugging the USB, and disconnecting the FLA from power. Similarly,
the PC may need to be rebooted.
\end{enumerate}
\item Open Cummins PowerSpec. 
\item Click on Advanced -> Settings and set the Adapter to the DPA4+
\item Click on Connect.
\item Click Read Data
\item If available:

\begin{enumerate}
\item Press the Fault Codes button and Save the report as a PDF file.
\item Press the Trip Information button and Save the report as a PDF file.
\item Press the Feature Settings button and Save the report as a PDF file.
\item Press the Sudden Decel button and Save the report as a PDF file.
\item Press the Dataplate button and Save the report as a PDF file.
\item Press the Duty Cycle button and Save the report as a PDF file.
\item Press the After Treatment button and Save the report as a PDF file.
\end{enumerate}
\item Exit passthrough mode by pressing a the Update (red) button on the
FLA. If an Internet Connection is available, then the data can be
uploaded and archived. Otherwise, shutdown.
\item Establish Internet connection with the FLA. An IP address other than
10.0.0.1 should show up on the FLA display. 
\item Scroll to the Upload to Server option on the FLA and upload the data
to the server.
\end{enumerate}

\section{\label{sec:Forensic-Write-Blocking}Forensic Write Blocking}

There are certain messages that should not be transmitted to the ECM
during a forensic examination. These messages are the kind that will
change the data stored on the ECM in a meaningful way. Resetting or
clearing fault codes, changing the ECM time, and resetting trip data
should always be avoided. 


\chapter{Suggested Practices}

This section is intended to provide a very brief overview of using
the FLA, for a more indebth guide, refer to the Field Guide.


\section{Connecting the FLA to the Internet}

In order to upload data packages to the FLA Porta, recieve FLA updates,
or update the FLA system time, the FLA needs to be connected to the
internet.
\begin{enumerate}
\item Plug the ethernet cable from the FLA into a known good wired Internet
connection that is on the same network as your computer.
\item Navigate to the FLA Title Screen \ref{enu:Upload-to-Server}. Line
2 should begin with ``IP:'' followed by 4 sets of numbers separated
by a ``.''
\item Navigate to the FLA Upload screen. If the green option is ``Upload'',
then the FLA has successfully connected to the internet. If not, refer
to the troubleshooting section.
\end{enumerate}

\section{Connect to the FLA Preview}

There are two different ways to connect to the FLA Preview, the in-office
way, and the in-field way. The in-office method requires the internet
connection outlined in the previous section, and is only required
if the operater needs the FLA Internet connection, i.e. to open a
support channel.


\subsection{In-Office Connection}
\begin{enumerate}
\item Connect the FLA to the internet as in the above section.
\item Using a web browser on a computer, navigate to the IP Address displayed
on the FLA. If you are unable to view the page, refer to the troubleshooting
section.
\end{enumerate}

\subsection{In-Field Connection}
\begin{enumerate}
\item If the computer being used is connected to a WiFi network, disconnect
from the network before proceeding.
\item Plug the Ethernet cable from the FLA to your computer's Ethernet port.
A USB to Ethernet adapter is available in the FLA kit if your computer
does not have an ethernet jack.
\item Navigate the FLA menu to the DHCP enable screen (Screen \ref{enu:DHCP})
and enable the DHCP Server on the FLA.

\begin{enumerate}
\item The IP Address of the FLA should now be 10.0.0.1
\end{enumerate}
\item Using a web browser on a computer, navigate to the IP Address displayed
on the FLA (10.0.0.1). If you are unable to view the page, refer to
the troubleshooting section.
\item When finished, do not forget to disable the DHCP service on the FLA.
Leaving the DHCP service on could cause problems if the FLA is plugged
into an office or home network.
\end{enumerate}

\section{Initiating FLA Support mode}

In order to provide the quickest and most complete support, it may
be necessary to place the FLA in support mode. This allows Synercon
Support staff to create a secure connection to an FLA anywhere with
an internet connection. It may also be refered to as a ``Support
Channel''.

To place the FLA in support mode, or ``Open a Support Channel'',
follow these steps
\begin{enumerate}
\item Verify a Synercon Support representative is ready for you to initiate
Support Mode.
\item Connect the FLA to the internet as described in section XXX.
\item Navigate to the FLA Preview site as described in section XXX.
\item Select the ``Support'' option from the menu at the top.
\item Click the ``Open the support channel'' button to initiate Support
Mode.
\item After a brief period, the page will change to say the connection was
successful, if the connection was not successful, refer to the troubleshooting
section.
\item Do not close the web page, or click the ``Back'' or ``Refresh''
buttons. It is important to stay on this page until the Synercon Support
representative is finished providing support.
\item When the Synercon Support representative begins working on your FLA,
the screen on the device will change to indicate a member is preforming
work.
\item When the Synercon Support representative is finished, the unit will
display a message indicating the work is finished.
\item The FLA will automatically reboot.
\item After the device loads the Title Screen, then click the ``Close Support''
button to return the FLA to normal operational mode.
\end{enumerate}

\section{Installing PC Software}


\subsection{RP 1210 Drivers}

The FLA can opperate as an RP1210 device. This is used for connecting
OEM software to ECMs. A computer does not need the RP 1210 drivers
to access the FLA Preview, and the FLA Portal. You will need administrative
privileges to install these drivers.

To install the RP1210 Drivers, follow these steps
\begin{enumerate}
\item Connect to the FLA Preview discribed in the above section (Ether In-Field,
or In-Office methods will work)
\item Select the Download Drivers option from the menu at the top of the
web page
\item Select ``Download Drivers''
\item Once the file is downloaded, extract the .zip file to a folder
\item Open the folder, and double-click ``DPA4Plus\_RP1210.exe''
\item Follow the on-screen instructions to install the drivers
\end{enumerate}
After the drivers are installed, they can be tested with the following
steps
\begin{enumerate}
\item Make sure the FLA is powered on, and connected to an ECM (It is possible
to test the drivers without a live ECM, to do so see the special instructions)

\begin{enumerate}
\item Skip steps 2,3,and 4 if not connected to a live ECM
\end{enumerate}
\item Begin a scan on the vehicle
\item After the scan has completed and you see the \textquotedbl{}Enable
RP1210 Mode\textquotedbl{} screen
\item Confirm the top line of the FLA read \textquotedbl{}CAN/J1939 Pass
Mode.\textquotedbl{}
\item Search on your computer for the program \textquotedbl{}Adapter Validation
Tool\textquotedbl{} and run it
\item Select \textquotedbl{}MD/HD\textquotedbl{}
\item Set the Vendor to \textquotedbl{}DG1210xx - Dearborn Group RP1210\textquotedbl{}
(the xx will likely be 32)
\item Set the Device to \textquotedbl{}DG DPA 4 Plus USB\textquotedbl{}

\begin{enumerate}
\item Skip to step 14 if the ECM does not support J1939
\end{enumerate}
\item Set Protocol to \textquotedbl{}J1939\textquotedbl{}
\item plug the USB cable into the FLA and the computer
\item Click \textquotedbl{}Run Test\textquotedbl{}
\item The text boxes labeled \textquotedbl{}RP1210 Status Window (PC-VDA)\textquotedbl{}
and \textquotedbl{}RP1210 Data Message Window (VDA-Vehicle)\textquotedbl{}
should turn green

\begin{enumerate}
\item If you are not connected to a live ECM, the \textquotedbl{}RP1210
Data Message Window\textquotedbl{} will turn red as there are no messages
\end{enumerate}
\item A text box will appear and will describe the results of the test.
\item Click \textquotedbl{}Quit\textquotedbl{} on the pop-up window titled
\textquotedbl{}Test Results Discussion\textquotedbl{}

\begin{enumerate}
\item If the ECM does not support J1708, the test is complete.If it does,
continue
\end{enumerate}
\item Change the protocol to J1708
\item On the FLA click the Green button to enable J1708 passthrough
\item Confirm the top line of the FLA read \textquotedbl{}J1708 Pass Mode
for\textquotedbl{}
\item Click \textquotedbl{}Run Test\textquotedbl{}
\item A text box will appear and will describe the results of the test
\item Click \textquotedbl{}Quit\textquotedbl{} on the pop-up window titled
\textquotedbl{}Test Results Discussion\textquotedbl{}
\item Close out of the adapter validation test
\end{enumerate}

\section{Preforming a download on a Heavy Vehicle}

This section provides a very high level overview of a typical download.
\begin{enumerate}
\item Plug in the FLA to the 9-Pin Diagnostics Connector
\item Turn on the ignition key to run
\item Run a vehicle Scan
\item Download Supported Data (if available)
\item Enable RP1210 Mode (Optional depending on the ECM and user's proceedure)

\begin{enumerate}
\item Connect the USB Cable
\item Run OEM Software
\end{enumerate}
\item Verify the data collected via the FLA Preview

\begin{enumerate}
\item If desired, repeat steps 3-6 until satisfied
\end{enumerate}
\item Shutdown the FLA
\item Reach a location that has a working Ethernet port
\item Connect the FLA to the Internet
\item Upload the data from the field
\item View the full report on the FLA Portal
\end{enumerate}

\section{Connecting to Engine Control Modules on a Bench}

When preforming bench downloads, it is recommended to use a Smart
Sensor Simulator to mimimize the chance of throwing additional faults
as compared to a traditional calibration harness. For more information
about Smart Sensor Simulators, visit \textbackslash{}url\{synercontechnologies.com\}.
Once connected, use the FLA in accordance to the recomended practices
outlined in the Field Guide and accompanying Quick Start guides.

The following are typical steps taken during a download.
\begin{enumerate}
\item Confirm the correct SSS is being used for the ECM
\item Connect the SSS to the ECM
\item Connect power to the SSS via the included power adapter
\item Press the Key Switch until the Key On lamp is iluminated, about two
seconds
\item Verify the Key On lamp stays iluminated for one second before proceeding.
If not, refer to the troubleshooting section
\item Plug in the FLA to the 9-Pin Diagnostics Connector located on the
SSS
\item Run a vehicle Scan
\item Download Supported Data (if available)
\item Enable RP1210 Mode (Optional depending on the ECM and user's proceedure)

\begin{enumerate}
\item Connect the USB Cable
\item Run OEM Software
\end{enumerate}
\item Verify the data collected via the FLA Preview

\begin{enumerate}
\item If desired, repeat steps 3-6 until satisfied
\item Skip to step 13 if there is a working Ethernet port where the bench
download is occuring
\end{enumerate}
\item Shutdown the FLA
\item Reach a location that has a working Ethernet port
\item Connect the FLA to the Internet
\item Upload the data from the field
\item View the full report on the FLA Portal
\end{enumerate}

\chapter{Troubleshooting}

This section is dedicated to helping operators troubleshoot common
issues with the FLA. Please consult this section and preform the suggested
steps before contacting Synercon Support with any problems. The steps
outlined in this section are intended to be a guide, express common
sense when preforming any action that may result in data-loss or damage
to the FLA.


\section{\label{sec:FailsToBoot}FLA Fails to Boot}

As the FLA boots, it may see something that it doesn't like and fails
to pass its own self test. Often restarting will fix the problem.
If connected to a truck, restart with the key in the off position
so only power is connected. 


\chapter{Definitions}

When dealing with heavy truck data, there are many acronyms, symbols
and phrases that are used to describe the data. This chapter presents
a list to help decipher some of the terms used for dealing with heavy
vehicle event data recorders.

\nomenclature{FLA}{The Forenisc Link Adapter (FLA) is a field rugged computer that is programmed to download data from heavy vehicles.}
\begin{description}
\item [{NIST}] National Institute for Standards and Technology (\href{http://www.nist.gov/}{http://www.nist.gov/})
\item [{RTC}] Real Time Clock
\end{description}

\chapter{Requirements}

In order to take full advantage of the Synercon products, such as
the FLA, FLA Preview, and FLA Portal, the user will need to have access
to additional software and ports such as those outlined in the following
sections.


\section{Supported Web Browsers}

In order for the FLA Preview and FLA Portal to work correctly, the
user will need to meet the following browser requirements.
\begin{description}
\item [{Internet\_Explorer}] at least version XXX (MM/YYYY)
\item [{Firefox}] at least version XXX (MM/YYYY)
\item [{Chrome}] at least version XXX (MM/YYYY)
\item [{Opera}] at least version XXX (MM/YYYY)
\end{description}

\section{Drivers}

In order to use the FLA in RP 1210 mode, the operator will need to
install the drivers first. This will require administrative privilages.
The drivers are provided with the FLA.


\section{OEM Software}

In order to extract data from ECMs that are not fully supported, the
operator will need the relative OEM software to aquire non-standards-based
data.


\chapter{Warnings}

Synercon Technologies does not warrante any products against user
caused problems. In order to prevent any damage to the FLA and data
collected by the FLA, always follow these guidelines.


\section{Powering the FLA}


\subsection*{Only use the included wall adapter or cigarette adapter to power
the FLA}

When not using a vehicle to power the FLA, only power the device using
the included power supplies. Never plug in a forign adapter, espeically
if the voltage exceeds 12 volts.


\subsection*{Never plug the FLA into a diagnostic port that has a vehicle voltage
above 26 volts}

When plugging in the FLA to a vehicle, ensure the voltage of the diagnostic
port does not exceed 26 volts. If in doubt, a multi-meter can be used
to verify voltage


\subsection*{Do not remove power from the FLA without first shutting down}

Unless specifically instructed to do so by a troubleshooting guide,
avoid removing power to the FLA without first shutting down. Use the
securing screws of the cables to secure all cables against accidental
dislodgement.


\subsection*{Do not expose the FLA to water}

The FLA is NOT waterproof, even powered off. Avoid dropping the unit
in puddles, or exposing the device to rainfall. Take care to ensure
a dry environment is avaliable before operating the FLA. If the unit
is exposed to water, even when powered off, the water can damage the
unit, and potentially any data stored on the unit. If exposed to water,
imediatly power the unit off and contact Synercon Support for instructions.


\subsection*{Do not expose to excess heat or cold}

Do not expose the FLA to temperatures exceeding XXXF, or below XXXF.


\chapter{Synercon Training}

Synercon Technologies provides classes to train users in the proper
use of FLAs in investigating vehicles. For more information about
classes, please contact XXX.


\chapter{Contact Synercon Techonlogies}

While we strive to provide the hightest quality products, if you have
an issue with our products and need to contact us, please send an
email to \textbackslash{}email\{support@synercontechnologies.com\}.

When requesting support, please include the following:

1) the product (i.e. FLA, or SSS)

2) the serial number (it is printed on the bottom of all Synercon
Products)

3) A detailed explanation of the issue

4) Any steps you have already taken to solve the issue

5) {[}optional{]} any screen shots, or pictures if you believe they
will help us understand the problem.

We will attempt to respond to you as quickly as possible.

\printindex{}
\end{document}
